\chapter{Capitolo 9 - Teoremi fondativi F01-F10: enunciati, proof sketch, conseguenze}

\section{1. Problema}

I capitoli precedenti hanno costruito il telaio teorico di OCT:

\begin{enumerate}
\def\labelenumi{\arabic{enumi}.}
\tightlist
\item
  crisi della sola validità sintattica;
\item
  passaggio ontologico da cosa a singolarità;
\item
  assiomatica O1-O7 della categoria ordinativa;
\item
  spazio di validità \texttt{S\_Omega} con regioni A/B/C e invarianti I1-I5.
\end{enumerate}

Questo non basta ancora per una fondazione completa. Una teoria non entra nella grammatica della scienza solo perché possiede un lessico elegante o una struttura assiomatica coerente. Diventa scienza quando produce enunciati teorematici che:

\begin{enumerate}
\def\labelenumi{\arabic{enumi}.}
\tightlist
\item
  sono formalmente tracciabili;
\item
  distinguono casi reali da casi apparenti;
\item
  guidano protocolli di prova e controprova.
\end{enumerate}

Il compito del capitolo è quindi fissare i \textbf{dieci teoremi fondativi F01-F10} in forma sufficientemente stabile da sostenere:

\begin{enumerate}
\def\labelenumi{\arabic{enumi}.}
\tightlist
\item
  il programma differenziale (Capitolo 10);
\item
  la critica dei limiti e dei controesempi (Capitolo 11);
\item
  la validazione empirico-formale dei capitoli metodologici successivi.
\end{enumerate}

La difficoltà è metodologica oltre che formale. Se i teoremi vengono espressi in forma troppo forte rispetto all'evidenza disponibile, la teoria degrada in overclaiming. Se sono espressi in forma troppo prudente, non generano avanzamento. Il capitolo adotta quindi una disciplina esplicita:

\begin{enumerate}
\def\labelenumi{\arabic{enumi}.}
\tightlist
\item
  enunciato formale;
\item
  ipotesi dichiarate;
\item
  proof sketch verificabile;
\item
  conseguenze teorico-operative;
\item
  stato di maturità (\texttt{validated}, \texttt{in\_proof}, \texttt{revise}).
\end{enumerate}

\subsection{1.1 Obiettivo del blocco fondativo}

L'obiettivo non è ``chiudere definitivamente'' l'intera teorematica OCT in un singolo passaggio, ma stabilire un blocco fondativo abbastanza forte da rendere non ridondante l'estensione ordinativa rispetto al classico.

In termini sintetici, i teoremi F01-F10 devono garantire tre risultati:

\begin{enumerate}
\def\labelenumi{\arabic{enumi}.}
\tightlist
\item
  \textbf{separazione}: esistono casi classically-valid ma non OCT-valid;
\item
  \textbf{chiusura condizionata}: la composizione ordinativa è governata da regole esplicite;
\item
  \textbf{conservatività}: il classico è recuperabile come caso limite.
\end{enumerate}

\subsection{1.2 Vincolo editoriale del capitolo}

Per coerenza con il template del libro, ogni teorema viene presentato in forma uniforme. Le prove complete saranno sviluppate nei cicli successivi di validazione. Qui fissiamo la versione canonica per uso nel manoscritto fondativo.

\begin{center}\rule{0.5\linewidth}{0.5pt}\end{center}

\section{2. Tesi del capitolo}

La tesi del capitolo è articolata in sei enunciati.

\begin{enumerate}
\def\labelenumi{\arabic{enumi}.}
\tightlist
\item
  I teoremi F01-F10 costituiscono un blocco fondativo coerente con O1-O7.
\item
  Il blocco dimostra che OCT non è una ridenominazione del classico, ma un'estensione conservativa selettiva.
\item
  La selettività ordinativa emerge da tre criteri congiunti: coerenza, emergenza, robustezza contestuale.
\item
  I teoremi fondativi producono conseguenze operative dirette su composizione, equivalenza, limiti, colimiti e osservazione.
\item
  Lo stato di maturità dei singoli teoremi è esplicito e tracciabile; nessun enunciato viene presentato come già chiuso quando non lo è.
\item
  Il blocco F01-F10 è sufficiente per sostenere il passaggio ai teoremi differenziali D01-D05 e applicativi A01-A03.
\end{enumerate}

\begin{center}\rule{0.5\linewidth}{0.5pt}\end{center}

\section{3. Definizioni canoniche}

Prima degli enunciati, fissiamo il lessico minimo del capitolo.

\subsection{Definizione 9.1 (Validità ordinativa)}

\texttt{OCT\_Valid\_Omega(D)} vale se e solo se:

\texttt{Syn(D)\ and\ Coh\_Omega(D)\ \textgreater{}=\ tau\ and\ Phi\_Omega(D)\ !=\ 0}.

\subsection{Definizione 9.2 (Equivalenza ordinativa)}

Due oggetti o strutture sono ordinativamente equivalenti in \texttt{Omega} se:

\begin{enumerate}
\def\labelenumi{\arabic{enumi}.}
\tightlist
\item
  sono correlabili da morfismi ammissibili in entrambe le direzioni;
\item
  preservano invarianti ordinativi rilevanti entro tolleranze dichiarate.
\end{enumerate}

\subsection{Definizione 9.3 (Composizione coerente)}

Una composizione \texttt{g\ o\ f} è coerente in \texttt{Omega} se:

\begin{enumerate}
\def\labelenumi{\arabic{enumi}.}
\tightlist
\item
  è sintatticamente definita;
\item
  soddisfa il gate ordinativo su \texttt{Coh} e \texttt{Phi};
\item
  resta stabile su vicinato contestuale ammissibile.
\end{enumerate}

\subsection{Definizione 9.4 (Controesempio ordinativo)}

Un controesempio ordinativo è un caso con:

\begin{enumerate}
\def\labelenumi{\arabic{enumi}.}
\tightlist
\item
  validità classica soddisfatta;
\item
  violazione almeno di una condizione ordinativa.
\end{enumerate}

\subsection{Definizione 9.5 (Recupero classico)}

Il recupero classico è la riduzione in cui i vincoli ordinativi vengono disattivati o resi trivialmente veri, mostrando che OCT contiene il quadro classico come caso limite.

\begin{center}\rule{0.5\linewidth}{0.5pt}\end{center}

\section{4. Sviluppo formale}

\subsection{4.1 Architettura dei teoremi F01-F10}

Il blocco fondativo è organizzato in quattro famiglie.

\begin{enumerate}
\def\labelenumi{\arabic{enumi}.}
\tightlist
\item
  \textbf{Differenza ontologica}: F01, F02, F07.
\item
  \textbf{Struttura compositiva}: F03, F04.
\item
  \textbf{Universalità selettiva}: F05, F06.
\item
  \textbf{Meta-struttura contestuale e riduzione}: F08, F09, F10.
\end{enumerate}

Questa architettura evita dispersione: ogni teorema copre un nodo critico del confronto classico/OCT.

\subsection{4.2 Teorema F01 - Non-collasso della singolarità}

\textbf{Teorema F01 (Non-collasso della singolarità).}\\
L'isomorfismo categoriale classico non implica, in generale, equivalenza ordinativa di funzione singolare in \texttt{Omega}.

Ipotesi:

\begin{enumerate}
\def\labelenumi{\arabic{enumi}.}
\tightlist
\item
  \texttt{A} e \texttt{B} sono isomorfi nel senso classico;
\item
  esiste singolarizzazione \texttt{S\_Omega} su \texttt{A} e \texttt{B};
\item
  almeno un invariante funzionale ordinativo differisce oltre soglia.
\end{enumerate}

Tesi:
\texttt{A\ \ \textasciitilde{}=\ \ B} non implica \texttt{A\ \textasciitilde{}ord,Omega\ B}.

Proof sketch:

\begin{enumerate}
\def\labelenumi{\arabic{enumi}.}
\tightlist
\item
  l'isomorfismo garantisce preservazione di struttura classica;
\item
  la funzione singolare può dipendere da profilo relazionale e traiettoria storica non catturati dall'isomorfismo;
\item
  se gli invarianti ordinativi divergono oltre tolleranza, l'equivalenza ordinativa decade.
\end{enumerate}

Conseguenze:

\begin{enumerate}
\def\labelenumi{\arabic{enumi}.}
\tightlist
\item
  blocca il collasso ``isomorfo = uguale in tutto'' nei domini reali complessi;
\item
  giustifica l'esistenza di un livello ordinativo distinto.
\end{enumerate}

Stato: \texttt{in\_proof}.

\subsection{4.3 Teorema F02 - Ammissibilità generativa dei morfismi}

\textbf{Teorema F02 (Ammissibilità generativa dei morfismi).}\\
Esistono morfismi formalmente ben tipizzati che non sono OCT-validi.

Ipotesi:

\begin{enumerate}
\def\labelenumi{\arabic{enumi}.}
\tightlist
\item
  \texttt{f:\ A-\textgreater{}B} è ben tipizzato;
\item
  \texttt{Coh\_Omega(f)} o \texttt{Phi\_Omega(f)} viola il gate.
\end{enumerate}

Tesi:
\texttt{Syn(f)} non implica \texttt{OCT\_Valid\_Omega(f)}.

Proof sketch:

\begin{enumerate}
\def\labelenumi{\arabic{enumi}.}
\tightlist
\item
  O1 assicura che il criterio ordinativo non contraddice la tipizzazione;
\item
  O2-O3 introducono vincoli aggiuntivi indipendenti dalla sola correttezza sintattica;
\item
  basta un caso con \texttt{Phi=0} o \texttt{Coh\textless{}tau} per ottenere il controesempio.
\end{enumerate}

Conseguenze:

\begin{enumerate}
\def\labelenumi{\arabic{enumi}.}
\tightlist
\item
  dimostra che la tipizzazione è necessaria ma non sufficiente;
\item
  fonda il concetto di filtro ordinativo.
\end{enumerate}

Stato: \texttt{validated} (in forma logica), con estensione empirica in corso.

\subsection{4.4 Teorema F03 - Chiusura compositiva coerente}

\textbf{Teorema F03 (Chiusura compositiva coerente).}\\
Sotto stabilità di coerenza locale, la composizione di morfismi OCT-validi resta OCT-valida.

Ipotesi:

\begin{enumerate}
\def\labelenumi{\arabic{enumi}.}
\tightlist
\item
  \texttt{f:\ A-\textgreater{}B}, \texttt{g:\ B-\textgreater{}C} sono OCT-validi;
\item
  il gate sulla composta \texttt{g\ o\ f} è soddisfatto;
\item
  le perturbazioni locali non spingono la composta sotto soglia.
\end{enumerate}

Tesi:
\texttt{OCT\_Valid\_Omega(f)\ and\ OCT\_Valid\_Omega(g)} implica \texttt{OCT\_Valid\_Omega(g\ o\ f)}.

Proof sketch:

\begin{enumerate}
\def\labelenumi{\arabic{enumi}.}
\tightlist
\item
  dalla composizione classica otteniamo definibilità sintattica;
\item
  dalla stabilità di coerenza locale otteniamo persistenza di \texttt{Coh\textgreater{}=tau};
\item
  dalla non-sterilità emergenziale otteniamo \texttt{Phi!=0} sulla composta.
\end{enumerate}

Conseguenze:

\begin{enumerate}
\def\labelenumi{\arabic{enumi}.}
\tightlist
\item
  rende possibile la costruzione di pipeline ordinative lunghe;
\item
  fornisce base per i risultati su sistemi modulari.
\end{enumerate}

Stato: \texttt{in\_proof}.

\subsection{4.5 Teorema F04 - Neutralità ordinativa dell'identità}

\textbf{Teorema F04 (Neutralità ordinativa dell'identità).}\\
\texttt{id\_A} è ordinativamente neutra solo se preserva la funzione singolare in \texttt{Omega}.

Ipotesi:

\begin{enumerate}
\def\labelenumi{\arabic{enumi}.}
\tightlist
\item
  \texttt{id\_A} è identità classica;
\item
  esiste misura ordinativa su stato prima/dopo identità.
\end{enumerate}

Tesi:
neutralità ordinativa richiede assenza di degradazione su \texttt{Coh} e \texttt{Phi}.

Proof sketch:

\begin{enumerate}
\def\labelenumi{\arabic{enumi}.}
\tightlist
\item
  la neutralità classica è tautologica sul piano sintattico;
\item
  sul piano ordinativo serve mostrare che l'auto-applicazione non induce perdita funzionale;
\item
  se si osserva decremento sistematico, l'identità non è neutra nel senso ordinativo.
\end{enumerate}

Conseguenze:

\begin{enumerate}
\def\labelenumi{\arabic{enumi}.}
\tightlist
\item
  raffina il concetto di identità per sistemi dinamici reali;
\item
  abilità diagnostica precoce su degenerazione interna.
\end{enumerate}

Stato: \texttt{in\_proof}.

\subsection{4.6 Teorema F05 - Universalità selettiva dei limiti}

\textbf{Teorema F05 (Universalità selettiva dei limiti).}\\
Non ogni limite classico è limite ordinativo reale.

Ipotesi:

\begin{enumerate}
\def\labelenumi{\arabic{enumi}.}
\tightlist
\item
  esistenza di limite classico \texttt{Lim(D)} per un diagramma \texttt{D};
\item
  disponibilità di misura ordinativa sul cono universale.
\end{enumerate}

Tesi:
\texttt{Lim(D)} è ordinativamente accettabile solo se soddisfa il gate \texttt{Coh/Phi}.

Proof sketch:

\begin{enumerate}
\def\labelenumi{\arabic{enumi}.}
\tightlist
\item
  la proprietà universale classica garantisce fattorizzazione strutturale;
\item
  la fattorizzazione può essere emergenzialmente sterile o coerentivamente fragile;
\item
  il filtro ordinativo seleziona solo limiti con funzione reale non degenerativa.
\end{enumerate}

Conseguenze:

\begin{enumerate}
\def\labelenumi{\arabic{enumi}.}
\tightlist
\item
  evita uso automatico di limiti formalmente corretti ma operativamente vuoti;
\item
  prepara la nozione di limite ordinativo nei capitoli 4-7.
\end{enumerate}

Stato: \texttt{revise}.

\subsection{4.7 Teorema F06 - Universalità selettiva dei colimiti}

\textbf{Teorema F06 (Universalità selettiva dei colimiti).}\\
Non ogni colimite classico preserva validità ordinativa.

Ipotesi:

\begin{enumerate}
\def\labelenumi{\arabic{enumi}.}
\tightlist
\item
  esistenza di \texttt{Colim(D)} in senso classico;
\item
  valutazione ordinativa dell'aggregazione risultante.
\end{enumerate}

Tesi:
\texttt{Colim(D)} è ordinativamente valido solo in presenza di coerenza sopra soglia ed emergenza non nulla.

Proof sketch:

\begin{enumerate}
\def\labelenumi{\arabic{enumi}.}
\tightlist
\item
  il colimite classico aggrega attraverso coconi universali;
\item
  alcune aggregazioni possono aumentare connessione ma annullare funzione emergente;
\item
  il criterio ordinativo distingue aggregazione strutturale da integrazione reale.
\end{enumerate}

Conseguenze:

\begin{enumerate}
\def\labelenumi{\arabic{enumi}.}
\tightlist
\item
  introduce selettività nell'uso di colimiti in sistemi complessi;
\item
  riduce falsi positivi di ``integrazione'' dovuti a sola connettività.
\end{enumerate}

Stato: \texttt{revise}.

\subsection{4.8 Teorema F07 - Equivalenza ordinativa vincolata}

\textbf{Teorema F07 (Equivalenza ordinativa vincolata).}\\
Equivalenza classica di categorie non implica equivalenza ordinativa.

Ipotesi:

\begin{enumerate}
\def\labelenumi{\arabic{enumi}.}
\tightlist
\item
  due categorie \texttt{C1}, \texttt{C2} sono equivalenti in senso classico;
\item
  esiste confronto di invarianti ordinativi su famiglie di oggetti/morfismi corrispondenti;
\item
  almeno un invariante chiave diverge oltre tolleranza.
\end{enumerate}

Tesi:
\texttt{C1\ \ \textasciitilde{}=\ \ C2} non implica \texttt{C1\ \ \textasciitilde{}=\ ord\ C2}.

Proof sketch:

\begin{enumerate}
\def\labelenumi{\arabic{enumi}.}
\tightlist
\item
  l'equivalenza classica preserva struttura fino a isomorfismo naturale;
\item
  gli invarianti ordinativi dipendono da stabilità contestuale e funzione emergente;
\item
  divergenze sistematiche su tali invarianti negano equivalenza ordinativa.
\end{enumerate}

Conseguenze:

\begin{enumerate}
\def\labelenumi{\arabic{enumi}.}
\tightlist
\item
  rafforza la non-ridondanza di OCT su scala categoriale globale;
\item
  giustifica benchmark comparativi tra architetture formalmente equivalenti.
\end{enumerate}

Stato: \texttt{in\_proof}.

\subsection{4.9 Teorema F08 - Non-degenerazione minimale}

\textbf{Teorema F08 (Non-degenerazione minimale).}\\
Un diagramma non vuoto con \texttt{Phi\_Omega(D)=0} non è OCT-reale.

Ipotesi:

\begin{enumerate}
\def\labelenumi{\arabic{enumi}.}
\tightlist
\item
  \texttt{D} è sintatticamente non vuoto e ben composto;
\item
  \texttt{Phi\_Omega(D)=0}.
\end{enumerate}

Tesi:
\texttt{OCT\_Valid\_Omega(D)} è falso.

Proof sketch:

\begin{enumerate}
\def\labelenumi{\arabic{enumi}.}
\tightlist
\item
  la definizione di validità ordinativa include \texttt{Phi!=0} come condizione necessaria;
\item
  la violazione di questa condizione nega direttamente l'ammissibilità;
\item
  dunque il diagramma è classificato in regione di sterilità.
\end{enumerate}

Conseguenze:

\begin{enumerate}
\def\labelenumi{\arabic{enumi}.}
\tightlist
\item
  stabilisce il criterio minimo anti-sterilità;
\item
  fornisce base per controesempi D02 e applicativo A03.
\end{enumerate}

Stato: \texttt{validated} (nella forma assiomatica minimale).

\subsection{4.10 Teorema F09 - Osservatore interno}

\textbf{Teorema F09 (Osservatore interno).}\\
La realtà ordinativa è invariabile per cambio di notazione equivalente, non per cambio arbitrario di contesto osservativo.

Ipotesi:

\begin{enumerate}
\def\labelenumi{\arabic{enumi}.}
\tightlist
\item
  trasformazione notazionale \texttt{N} che preserva struttura e protocollo;
\item
  trasformazione contestuale \texttt{T} che altera variabili osservative sostanziali.
\end{enumerate}

Tesi:

\begin{enumerate}
\def\labelenumi{\arabic{enumi}.}
\tightlist
\item
  \texttt{N} preserva classificazione ordinativa;
\item
  \texttt{T} può alterarla in modo legittimo.
\end{enumerate}

Proof sketch:

\begin{enumerate}
\def\labelenumi{\arabic{enumi}.}
\tightlist
\item
  l'invarianza sintattica garantisce equivalenza descrittiva;
\item
  il contesto osservativo entra nella definizione stessa di \texttt{Coh} e \texttt{Phi};
\item
  cambiare \texttt{Omega} senza regole di trasporto equivale a cambiare il problema.
\end{enumerate}

Conseguenze:

\begin{enumerate}
\def\labelenumi{\arabic{enumi}.}
\tightlist
\item
  chiarisce il ruolo epistemico del contesto;
\item
  impedisce l'uso improprio di confronti cross-domain non normalizzati.
\end{enumerate}

Stato: \texttt{in\_proof}.

\subsection{4.11 Teorema F10 - Recupero classico}

\textbf{Teorema F10 (Recupero classico).}\\
La category theory classica si recupera come caso limite di OCT quando i vincoli ordinativi sono disattivati.

Ipotesi:

\begin{enumerate}
\def\labelenumi{\arabic{enumi}.}
\tightlist
\item
  struttura OCT con componente classica intatta;
\item
  vincoli ordinativi resi triviali o inattivi (\texttt{tau} minimale, gating non restrittivo in regime formale astratto).
\end{enumerate}

Tesi:
il comportamento strutturale risultante coincide con il quadro classico.

Proof sketch:

\begin{enumerate}
\def\labelenumi{\arabic{enumi}.}
\tightlist
\item
  O1 garantisce conservazione sintattica;
\item
  rimuovendo il potere selettivo di O2-O7 resta la sola grammatica classica;
\item
  ne segue riduzione conservativa.
\end{enumerate}

Conseguenze:

\begin{enumerate}
\def\labelenumi{\arabic{enumi}.}
\tightlist
\item
  esclude la lettura di OCT come rottura antitetica;
\item
  stabilisce continuità teorica e compatibilità didattica.
\end{enumerate}

Stato: \texttt{in\_proof}.

\subsection{4.12 Coerenza interna del blocco F01-F10}

Il blocco presenta una progressione logica controllata:

\begin{enumerate}
\def\labelenumi{\arabic{enumi}.}
\tightlist
\item
  F01-F02 dimostrano differenza locale tra validità classica e validità ordinativa;
\item
  F03-F04 governano composizione e identità in chiave non-degenerativa;
\item
  F05-F06 estendono la selettività a limiti e colimiti;
\item
  F07 scala la differenza al livello di equivalenza categoriale;
\item
  F08 fornisce criterio minimo anti-sterilità;
\item
  F09 controlla il ruolo dell'osservazione;
\item
  F10 chiude con riduzione conservativa al classico.
\end{enumerate}

Questa sequenza è intenzionale: evita sia frammentazione teorematica sia sovrapposizione ridondante.

\subsection{4.13 Matrice di maturità teorematica}

Per rendere auditabile lo stato attuale:

\begin{enumerate}
\def\labelenumi{\arabic{enumi}.}
\tightlist
\item
  \texttt{validated}: F02 (forma logica), F08 (forma assiomatica minima);
\item
  \texttt{in\_proof}: F01, F03, F04, F07, F09, F10;
\item
  \texttt{revise}: F05, F06.
\end{enumerate}

Questa distribuzione è coerente con la strategia editoriale del libro: dichiarare chiaramente dove la teoria è già forte e dove richiede ulteriore chiusura.

\subsection{4.14 Mappa dipendenze e ordine logico di dimostrazione}

Per evitare salti inferenziali, riportiamo una mappa di dipendenze operative, utile sia per la lettura sia per la validazione.

\begin{enumerate}
\def\labelenumi{\arabic{enumi}.}
\tightlist
\item
  \textbf{F01} dipende da O1, M01, M05: apre la separazione isomorfismo/funzione singolare.
\item
  \textbf{F02} dipende da O2, M02: introduce la frattura tipizzazione/ammissibilità.
\item
  \textbf{F03} dipende da O3, M03: formalizza la chiusura compositiva condizionata.
\item
  \textbf{F04} dipende da O1, O3, M04: raffina la nozione di identità in senso dinamico.
\item
  \textbf{F05} dipende da O3, O5, M08: seleziona limiti classici tramite gate ordinativo.
\item
  \textbf{F06} dipende da O3, O5, M09: seleziona colimiti tramite controllo di emergenza.
\item
  \textbf{F07} dipende da O1, O5, O7, M06: scala la differenza al livello di equivalenza categoriale.
\item
  \textbf{F08} dipende da O6, M07, M27: formalizza la non-degenerazione minima.
\item
  \textbf{F09} dipende da O7, M23, M24: distingue invarianza notazionale da invarianza contestuale.
\item
  \textbf{F10} dipende da O1-O7, M30: chiude con riduzione conservativa al classico.
\end{enumerate}

Da questa lista emerge un ordine tecnico consigliato per la chiusura formale:

\begin{enumerate}
\def\labelenumi{\arabic{enumi}.}
\tightlist
\item
  chiudere prima F03/F08/F10 (nucleo di coerenza globale);
\item
  consolidare F01/F02/F07 (dimostrazione di non-ridondanza);
\item
  rifinire F04/F05/F06/F09 (estensioni strutturali e contestuali).
\end{enumerate}

Razionale:

\begin{enumerate}
\def\labelenumi{\arabic{enumi}.}
\tightlist
\item
  senza F03 e F08 la nozione di ammissibilità resta fragile;
\item
  senza F10 la conservatività resta implicita;
\item
  senza F01/F07 la differenza col classico rischia di apparire solo retorica.
\end{enumerate}

Questa gerarchia non è un vincolo assoluto, ma una strategia che minimizza lavoro ridondante nei cicli di validazione.

\subsection{4.15 Schemi di confutazione per F01-F10}

Per mantenere il capitolo falsificabile, ogni teorema deve avere almeno un criterio di confutazione esplicito.

\begin{enumerate}
\def\labelenumi{\arabic{enumi}.}
\tightlist
\item
  \textbf{F01 fallisce} se si dimostra che ogni isomorfismo classico preserva sempre tutti gli invarianti ordinativi rilevanti.
\item
  \textbf{F02 fallisce} se non esiste alcun morfismo ben tipizzato che violi gate \texttt{Coh/Phi}.
\item
  \textbf{F03 fallisce} se si produce una composta che soddisfa ipotesi e tuttavia perde validità ordinativa.
\item
  \textbf{F04 fallisce} se un'identità con perdita misurabile su \texttt{Phi} resta comunque neutra per definizione formale adottata.
\item
  \textbf{F05 fallisce} se ogni limite classico risulta sempre ordinativamente valido senza eccezioni.
\item
  \textbf{F06 fallisce} se ogni colimite classico preserva sempre coerenza ed emergenza sopra soglia.
\item
  \textbf{F07 fallisce} se ogni equivalenza classica di categorie induce sempre equivalenza ordinativa.
\item
  \textbf{F08 fallisce} se esiste un diagramma con \texttt{Phi=0} classificabile comunque come OCT-reale.
\item
  \textbf{F09 fallisce} se il cambio arbitrario di contesto non altera mai classificazione ordinativa.
\item
  \textbf{F10 fallisce} se la disattivazione dei vincoli ordinativi non restituisce il comportamento classico.
\end{enumerate}

Questi criteri hanno una funzione precisa:

\begin{enumerate}
\def\labelenumi{\arabic{enumi}.}
\tightlist
\item
  impedire che i teoremi diventino enunciati non falsificabili;
\item
  guidare la progettazione dei benchmark in modo mirato;
\item
  rendere trasparente il passaggio \texttt{in\_proof\ -\textgreater{}\ validated}.
\end{enumerate}

In termini di metodo, OCT non richiede che ogni criterio di confutazione sia già osservato empiricamente, ma richiede che sia dichiarabile in forma controllata. La teoria resta così aperta alla revisione senza perdere struttura.

\begin{center}\rule{0.5\linewidth}{0.5pt}\end{center}

\section{5. Proposizioni e teoremi locali}

\subsection{Proposizione 9.1 (Copertura del nucleo assiomatico)}

Enunciato:
Ogni assioma O1-O7 è utilizzato da almeno un teorema del blocco F01-F10.

Intuizione:
Il programma fondativo non lascia assiomi ``decorativi'': ogni assioma produce conseguenze.

Stato: \texttt{validated}.

\subsection{Proposizione 9.2 (Non-ridondanza locale minima)}

Enunciato:
Esiste almeno una coppia di teoremi \texttt{Fi}, \texttt{Fj} con classi di conseguenze non equivalenti.

Intuizione:
Per esempio, F03 governa chiusura compositiva, F09 governa invarianza osservativa: non sono riducibili l'uno all'altro.

Stato: \texttt{validated}.

\subsection{Teorema 9.3 (Sufficienza strutturale del blocco fondativo)}

Ipotesi:

\begin{enumerate}
\def\labelenumi{\arabic{enumi}.}
\tightlist
\item
  validità di F01-F04 e F08-F10 in forma operativa;
\item
  revisione controllata di F05-F06.
\end{enumerate}

Tesi:
Il blocco è sufficiente per fondare il passaggio ai teoremi differenziali D01-D05 senza incoerenze strutturali.

Proof sketch:

\begin{enumerate}
\def\labelenumi{\arabic{enumi}.}
\tightlist
\item
  la differenza classico/OCT è già dimostrata localmente e globalmente;
\item
  la composizione e l'identità ordinative sono regolate;
\item
  la dimensione contestuale è formalmente inclusa;
\item
  il recupero classico impedisce contraddizioni di base.
\end{enumerate}

Conseguenze:
fornisce la base per il Capitolo 10.

Stato: \texttt{in\_proof}.

\subsection{Teorema 9.4 (Compatibilità con lo spazio di validità del Capitolo 8)}

Ipotesi:

\begin{enumerate}
\def\labelenumi{\arabic{enumi}.}
\tightlist
\item
  adozione di regioni A/B/C e invarianti I1-I5;
\item
  interpretazione dei teoremi F01-F10 tramite \texttt{Coh}, \texttt{Phi}, \texttt{Delta}.
\end{enumerate}

Tesi:
Ogni teorema fondativo ammette una traduzione in criteri classificativi nello spazio \texttt{S\_Omega}.

Proof sketch:

\begin{enumerate}
\def\labelenumi{\arabic{enumi}.}
\tightlist
\item
  F02/F08 mappano direttamente alla separazione tra A e B/C;
\item
  F03/F04 usano dinamica e persistenza in A;
\item
  F05-F07 distinguono universalità/ equivalenza con gate ordinativo;
\item
  F09-F10 governano trasporto contestuale e riduzione limite.
\end{enumerate}

Conseguenze:
chiude il circuito assiomi-metriche-teoremi.

Stato: \texttt{validated} (a livello di schema).

\subsection{Teorema 9.5 (Tracciabilità necessaria della maturità teorematica)}

Ipotesi:

\begin{enumerate}
\def\labelenumi{\arabic{enumi}.}
\tightlist
\item
  per ogni teorema \texttt{Fxx} è disponibile stato (\texttt{validated}, \texttt{in\_proof}, \texttt{revise});
\item
  esiste claim ledger con mappatura evidenza \texttt{S0-S3};
\item
  ogni variazione di stato è accompagnata da motivazione e riferimento a nuove prove o controprove.
\end{enumerate}

Tesi:
La maturità del blocco F01-F10 è auditabile in modo inter-soggettivo.

Proof sketch:

\begin{enumerate}
\def\labelenumi{\arabic{enumi}.}
\tightlist
\item
  la tracciabilità degli stati rende esplicito il grado di fiducia corrente;
\item
  la mappatura \texttt{S0-S3} lega forma teorica e qualità dell'evidenza;
\item
  la motivazione delle transizioni impedisce salti non giustificati;
\item
  segue che la maturità non dipende da autorità implicita, ma da log verificabile.
\end{enumerate}

Conseguenze:

\begin{enumerate}
\def\labelenumi{\arabic{enumi}.}
\tightlist
\item
  migliora la qualità della revisione tra pari;
\item
  riduce il rischio di regressioni silenziose nelle versioni successive;
\item
  favorisce integrazione con pipeline pubbliche (GitHub/Zenodo/OSF/HF) senza perdita di contesto.
\end{enumerate}

Stato: \texttt{validated}.

\begin{center}\rule{0.5\linewidth}{0.5pt}\end{center}

\section{6. Implicazioni operative}

\subsection{6.1 Per la stesura dei capitoli 10 e 11}

Il blocco F01-F10 fornisce una traccia precisa:

\begin{enumerate}
\def\labelenumi{\arabic{enumi}.}
\tightlist
\item
  i teoremi differenziali dovranno mostrare scarti strutturali misurabili;
\item
  i controesempi del Capitolo 11 dovranno riferirsi esplicitamente a una violazione di ipotesi Fxx.
\end{enumerate}

In questo modo, la narrativa resta controllata da logica teorematica.

\subsection{6.2 Per la validazione empirico-formale}

Ogni teorema con stato \texttt{in\_proof} o \texttt{revise} dovrebbe avere:

\begin{enumerate}
\def\labelenumi{\arabic{enumi}.}
\tightlist
\item
  un set minimo di casi positivi;
\item
  un set minimo di controesempi;
\item
  criterio esplicito di fallimento.
\end{enumerate}

Questa disciplina accelera il passaggio da manoscritto a preprint tecnico robusto.

\subsection{6.3 Per la governance delle release}

Per ogni release del framework OCT è consigliato allegare:

\begin{enumerate}
\def\labelenumi{\arabic{enumi}.}
\tightlist
\item
  tabella F01-F10 con stato aggiornato;
\item
  modifiche di ipotesi o soglie;
\item
  evidenze nuove che motivano il cambio di stato.
\end{enumerate}

Questo rende trasparente l'evoluzione teorica e riduce ambiguità tra versioni.

\subsection{6.4 Per i benchmark su repository pubblici}

La pubblicazione su GitHub/Zenodo/OSF/HF beneficia di una convenzione uniforme:

\begin{enumerate}
\def\labelenumi{\arabic{enumi}.}
\tightlist
\item
  benchmark etichettati con riferimento al teorema target (\texttt{F03}, \texttt{F08}, ecc.);
\item
  output classificato in termini A/B/C/D;
\item
  report con claim ledger coerente con \texttt{S0-S3}.
\end{enumerate}

Così il lettore esterno collega immediatamente risultato empirico e nodo teorico.

\subsection{6.5 Per la comunicazione interdisciplinare}

Il blocco F01-F10 è utile anche oltre la matematica pura, perché:

\begin{enumerate}
\def\labelenumi{\arabic{enumi}.}
\tightlist
\item
  traduce differenze ontologiche in criteri operativi;
\item
  permette dialogo con scienze computazionali, sociali e cognitive senza perdere formalità.
\end{enumerate}

La condizione è mantenere sempre separati:

\begin{enumerate}
\def\labelenumi{\arabic{enumi}.}
\tightlist
\item
  enunciato formale;
\item
  interpretazione dominio-specifica;
\item
  evidenza empirica.
\end{enumerate}

\subsection{6.6 Piano di chiusura operativo per F01-F10}

Per portare il blocco verso stato preprint-ready è utile una procedura in tre cicli.

\textbf{Ciclo A - Chiusura nucleo}

\begin{enumerate}
\def\labelenumi{\arabic{enumi}.}
\tightlist
\item
  chiusura formale di F03, F08, F10;
\item
  definizione univoca dei modelli di perturbazione per robustezza locale;
\item
  validazione incrociata su almeno due famiglie di benchmark indipendenti.
\end{enumerate}

\textbf{Ciclo B - Chiusura differenziale}

\begin{enumerate}
\def\labelenumi{\arabic{enumi}.}
\tightlist
\item
  consolidamento F01, F02, F07 con controesempi costruttivi tracciati;
\item
  misurazione esplicita della divergenza tra equivalenza classica e ordinativa;
\item
  report comparativo con matrix A/B/C/D.
\end{enumerate}

\textbf{Ciclo C - Chiusura universale e contestuale}

\begin{enumerate}
\def\labelenumi{\arabic{enumi}.}
\tightlist
\item
  revisione tecnica di F04, F05, F06, F09;
\item
  stress test su limiti/colimiti e trasporto tra contesti;
\item
  decisione finale di stato (\texttt{validated} o \texttt{revise}) con giustificazione pubblica.
\end{enumerate}

Questa sequenza è pensata per massimizzare affidabilità e minimizzare conflitti interpretativi.

\subsection{\texorpdfstring{6.7 Criteri minimi di accettazione per passaggio a \texttt{validated}}{6.7 Criteri minimi di accettazione per passaggio a validated}}

Per uniformità editoriale, proponiamo che un teorema Fxx passi a \texttt{validated} solo se rispetta cinque condizioni.

\begin{enumerate}
\def\labelenumi{\arabic{enumi}.}
\tightlist
\item
  enunciato in forma tipizzata e non ambigua;
\item
  proof sketch riproducibile da revisore indipendente;
\item
  almeno un controesempio limite discusso;
\item
  relazione esplicita con almeno un assioma O1-O7 e con lo spazio \texttt{S\_Omega};
\item
  aggiornamento coerente di claim ledger e changelog teorico.
\end{enumerate}

Il vantaggio è doppio:

\begin{enumerate}
\def\labelenumi{\arabic{enumi}.}
\tightlist
\item
  qualità formale più alta;
\item
  governance più trasparente su release successive del manoscritto.
\end{enumerate}

\subsection{6.8 Indicatore sintetico di avanzamento del blocco fondativo}

Per monitorare il progresso in modo semplice, proponiamo un indicatore editoriale:

\texttt{Prog\_F\ :=\ (N\_validated\ +\ 0.5*N\_in\_proof)\ /\ 10}.

Interpretazione:

\begin{enumerate}
\def\labelenumi{\arabic{enumi}.}
\tightlist
\item
  \texttt{Prog\_F=1} indica blocco fondativo completamente chiuso;
\item
  valori intermedi indicano avanzamento parziale con peso ridotto per i teoremi non ancora validati;
\item
  il valore non sostituisce il giudizio scientifico, ma aiuta a comunicare lo stato del progetto a revisori e collaboratori.
\end{enumerate}

Nello stato corrente, l'indicatore va letto insieme alla distribuzione degli stati per evitare semplificazioni: un progresso numerico alto con teoremi critici ancora \texttt{revise} richiede comunque prudenza editoriale.

\begin{center}\rule{0.5\linewidth}{0.5pt}\end{center}

\section{7. Limiti e non-claim}

Limiti espliciti del capitolo:

\begin{enumerate}
\def\labelenumi{\arabic{enumi}.}
\tightlist
\item
  molte prove sono in stato di proof sketch e non ancora chiusura completa;
\item
  F05 e F06 richiedono ulteriore raffinamento tecnico sui casi limite;
\item
  la parte quantitativa di robustezza contestuale richiede specifiche dominio-dipendenti;
\item
  la matrice di maturità non sostituisce validazione indipendente.
\end{enumerate}

Failure mode riconosciuti:

\begin{enumerate}
\def\labelenumi{\arabic{enumi}.}
\tightlist
\item
  confondere \texttt{in\_proof} con \texttt{validated};
\item
  usare F10 (recupero classico) per minimizzare il valore selettivo di F01-F09;
\item
  applicare F09 senza protocollo di trasporto contestuale;
\item
  dichiarare equivalenze ordinative senza controllo invarianti.
\item
  promuovere a \texttt{validated} teoremi che non hanno ancora criteri di confutazione espliciti.
\end{enumerate}

Box C - Non-claim

Questo capitolo non implica:

\begin{enumerate}
\def\labelenumi{\arabic{enumi}.}
\tightlist
\item
  che tutti i teoremi F01-F10 siano già definitivamente dimostrati;
\item
  che ogni dominio richieda lo stesso set di soglie e normalizzazioni;
\item
  che la sola formalizzazione sostituisca il ciclo di benchmark e revisione indipendente.
\end{enumerate}

\begin{center}\rule{0.5\linewidth}{0.5pt}\end{center}

\section{8. Claim Ledger (S0-S3)}

{\def\LTcaptype{none} % do not increment counter
\begin{longtable}[]{@{}lllll@{}}
\toprule\noalign{}
ID & Claim & Tipo & Evidenza & Stato \\
\midrule\noalign{}
\endhead
\bottomrule\noalign{}
\endlastfoot
C9.1 & Il blocco F01-F10 è coerente con il kernel O1-O7 & formale & S2 & in\_review \\
C9.2 & Esistono casi classically-valid ma non OCT-valid (F02/F08) & formale & S1 & validated \\
C9.3 & L'isomorfismo classico non implica equivalenza ordinativa (F01) & formale & S2 & in\_proof \\
C9.4 & La chiusura compositiva ordinativa è condizionata, non automatica (F03) & formale & S2 & in\_proof \\
C9.5 & Limiti e colimiti classici richiedono filtro selettivo ordinativo (F05/F06) & formale & S2 & revise \\
C9.6 & L'equivalenza categoriale classica non garantisce equivalenza ordinativa (F07) & formale & S2 & in\_proof \\
C9.7 & La realtà ordinativa dipende da \texttt{Omega} in modo formalmente controllato (F09) & epistemico-formale & S2 & in\_proof \\
C9.8 & Il quadro classico è recuperabile come caso limite (F10) & metateorico & S2 & in\_proof \\
C9.9 & Il blocco F01-F10 è sufficiente ad avviare il programma D01-D05 & programmatico & S2 & in\_review \\
\end{longtable}
}

\begin{center}\rule{0.5\linewidth}{0.5pt}\end{center}

\section{9. Bridge al Capitolo 10}

Il Capitolo 9 ha fissato il cuore teorematico fondativo:

\begin{enumerate}
\def\labelenumi{\arabic{enumi}.}
\tightlist
\item
  differenza formale tra validità classica e validità ordinativa;
\item
  regole su composizione, identità, universalità selettiva;
\item
  ruolo del contesto osservativo e recupero classico.
\end{enumerate}

Il Capitolo 10 svilupperà i teoremi differenziali D01-D05 e applicativi A01-A03 come test di stress del blocco fondativo, con attenzione a:

\begin{enumerate}
\def\labelenumi{\arabic{enumi}.}
\tightlist
\item
  casi di distinzione forte tra strutture formalmente simili;
\item
  perdita ontologica e dualità condizionata;
\item
  impatto su pipeline IA e sistemi narrativi complessi.
\end{enumerate}

In sintesi: se il Capitolo 9 stabilisce le leggi di base, il Capitolo 10 ne misura la capacità discriminante in scenari concreti e comparativi.

\begin{center}\rule{0.5\linewidth}{0.5pt}\end{center}

\section{Riferimenti}

Riferimenti di framework interno:

\begin{enumerate}
\def\labelenumi{\arabic{enumi}.}
\tightlist
\item
  Ghioni, F. (2026). \texttt{TE\_CORE\_v5.1.md}.
\item
  Ghioni, F. (2026). \texttt{Ordinative\_Set\_Theory\_OST\_A\_Concise\_Guide\_For\_AI\_v2\_1.md}.
\item
  \texttt{OCT\_THEOREM\_PROGRAM\_v0\_1.md}.
\item
  \texttt{OCT\_TYPED\_FORMAL\_SPEC\_v0\_1.md}.
\item
  \texttt{OCT\_CLASSICAL\_TO\_ORDINATIVE\_MAP\_v0\_1.md}.
\item
  \texttt{OCT\_FOUNDATIONAL\_BOOK\_CHAPTER\_08\_v1\_0.md}.
\end{enumerate}

Riferimenti primari:

\begin{enumerate}
\def\labelenumi{\arabic{enumi}.}
\tightlist
\item
  Eilenberg, S., Mac Lane, S. (1945). \emph{General Theory of Natural Equivalences}.
\item
  Mac Lane, S. (1998). \emph{Categories for the Working Mathematician}.
\item
  Riehl, E. (2016). \emph{Category Theory in Context}.
\item
  Spivak, D. (2014). \emph{Category Theory for the Sciences}.
\end{enumerate}
